\documentclass[a4paper,11pt]{ltjsarticle}
\usepackage{luatexja}
\usepackage{amsmath,amssymb,amsthm}
\usepackage{longtable,booktabs,array,calc}
\usepackage{hyperref}
\usepackage[margin=25mm]{geometry}

\newtheorem{theorem}{定理}[section]
\newtheorem{proposition}[theorem]{命題}
\newtheorem{lemma}[theorem]{補題}
\newtheorem{corollary}[theorem]{系}
\newtheorem{definition}[theorem]{定義}
\newtheorem{remark}[theorem]{注意}
\newtheorem{observation}[theorem]{観察}

\providecommand{\tightlist}{\setlength{\itemsep}{0pt}\setlength{\parskip}{0pt}}
\newcounter{none}

\begin{document}

\section{sine-Gordon方程式 ──
位相的ソリトンの基礎理論}\label{sine-gordonux65b9ux7a0bux5f0f-ux4f4dux76f8ux7684ux30bdux30eaux30c8ux30f3ux306eux57faux790eux7406ux8ad6}

\subsection{── 非線形波動方程式における位相的解の分類と衝突特性
──}\label{ux975eux7ddaux5f62ux6ce2ux52d5ux65b9ux7a0bux5f0fux306bux304aux3051ux308bux4f4dux76f8ux7684ux89e3ux306eux5206ux985eux3068ux885dux7a81ux7279ux6027}

\textbf{著者:} 木原 範昭 (WF System Co., Ltd.)

\textbf{日付:} 2026年4月

\textbf{DOI:} 10.5281/zenodo.19650966

\begin{center}\rule{0.5\linewidth}{0.5pt}\end{center}

\subsection{0. 本稿の目的}\label{ux672cux7a3fux306eux76eeux7684}

本稿は、非線形波動方程式の一つである \textbf{sine-Gordon方程式}
の数理的構造を概観する。sine-Gordon方程式は、位相的ソリトン（トポロジカル・ソリトン）を解として持つ最も基本的な非線形方程式であり、素粒子物理学、凝縮系物理学、場の理論において広く応用されている。

本稿では、sine-Gordon方程式の定義、解の分類（キンク、反キンク、ブリーザー）、衝突時の位相シフト、および物理的応用を、原論文に遡って整理する。さらに、ソリトン解が有限個のパラメータで厳密に特徴づけられること、および広がりパラメータ
\(\sigma \to 0\)
の極限（矩形波極限）においてソリトンの全ての構造的性質が保持されることを証明する。加えて、sine-Gordon理論の適用限界を明示し（§9.8）、これを克服する代替モデルとして、異なる比重を持つ完全流体の任意位相における完全弾性衝突の波動方程式による定式化を提示する（§11）。

\begin{center}\rule{0.5\linewidth}{0.5pt}\end{center}

\subsection{1.
sine-Gordon方程式の定義}\label{sine-gordonux65b9ux7a0bux5f0fux306eux5b9aux7fa9}

\subsubsection{1.1 方程式}\label{ux65b9ux7a0bux5f0f}

1+1次元の sine-Gordon方程式は以下で定義される {[}1, 2{]}。

\[\frac{\partial^2 \phi}{\partial t^2} - \frac{\partial^2 \phi}{\partial x^2} + \sin(\phi) = 0\]

ここで \(\phi = \phi(x, t)\) はスカラー場（位相場）、\(x\)
は空間座標、\(t\) は時間座標である。

自然単位系（\(c = 1\), \(m = 1\)）を用いている。一般の質量パラメータ
\(m\) と光速 \(c\) を復元すると、

\[\frac{1}{c^2}\frac{\partial^2 \phi}{\partial t^2} - \frac{\partial^2 \phi}{\partial x^2} + \frac{m^2 c^2}{\hbar^2}\sin(\phi) = 0\]

となる。

\subsubsection{1.2 名前の由来}\label{ux540dux524dux306eux7531ux6765}

sine-Gordon方程式の名称は、Klein-Gordon方程式

\[\frac{\partial^2 \phi}{\partial t^2} - \frac{\partial^2 \phi}{\partial x^2} + m^2 \phi = 0\]

との類似に由来する。線形項 \(m^2 \phi\) を非線形項 \(\sin(\phi)\)
に置き換えたものであり、``sine'' と ``Gordon''
を組み合わせた語呂合わせである。この名称は Rubinstein {[}3{]}
によって導入された。

\subsubsection{1.3
ラグランジアン密度}\label{ux30e9ux30b0ux30e9ux30f3ux30b8ux30a2ux30f3ux5bc6ux5ea6}

sine-Gordon方程式は以下のラグランジアン密度から導かれる {[}4{]}。

\[\mathcal{L} = \frac{1}{2}\left(\frac{\partial \phi}{\partial t}\right)^2 - \frac{1}{2}\left(\frac{\partial \phi}{\partial x}\right)^2 - (1 - \cos\phi)\]

ポテンシャル \(V(\phi) = 1 - \cos\phi\) は
\(\phi = 2n\pi\)（\(n \in \mathbb{Z}\)）で最小値0を取る。この無限個の縮退した真空の存在が、位相的ソリトンの起源である。

\subsubsection{1.4 歴史的背景}\label{ux6b74ux53f2ux7684ux80ccux666f}

sine-Gordon方程式は、以下の文脈で独立に導入された。

\begin{itemize}
\tightlist
\item
  \textbf{結晶転位理論}（1939年）：Frenkel と Kontorova
  が結晶格子中の転位の運動を記述するモデルとして導入した {[}1{]}。
\item
  \textbf{微分幾何学}（19世紀）：負の定曲率を持つ曲面（擬球面）の理論において、Bour
  {[}5{]} および Enneper {[}6{]} が同型の方程式を導出していた。Bianchi
  {[}7{]} はこの方程式のBäcklund変換を研究した。
\item
  \textbf{素粒子物理学}（1960年代）：Skyrme
  がバリオンの非線形場理論モデルとして sine-Gordon型の方程式を用いた
  {[}8{]}。Perring と Skyrme
  はキンク-反キンク衝突の数値計算を最初に行った {[}9{]}。
\end{itemize}

\begin{center}\rule{0.5\linewidth}{0.5pt}\end{center}

\subsection{2. 解の分類}\label{ux89e3ux306eux5206ux985e}

\subsubsection{\texorpdfstring{2.1 キンク解（巻き数
\(+1\)）}{2.1 キンク解（巻き数 +1）}}\label{ux30adux30f3ux30afux89e3ux5dfbux304dux6570-1}

キンク解は、\(\phi\) が \(0\) から \(2\pi\)
へ滑らかに遷移する局在波である {[}2, 9{]}。

\[\phi_K(x, t) = 4\arctan\left(\exp\left(\frac{x - vt}{\sqrt{1 - v^2}}\right)\right)\]

ここで \(v\)（\(|v| < 1\)）はキンクの速度である。

\textbf{性質:}

\begin{itemize}
\tightlist
\item
  \(x \to -\infty\) で \(\phi \to 0\)、\(x \to +\infty\) で
  \(\phi \to 2\pi\)
\item
  位相が \(2\pi\) だけ増加する ── \textbf{巻き数（トポロジカル荷）}
  \(N = +1\)
\item
  分母 \(\sqrt{1 - v^2}\) はローレンツ因子であり、特殊相対論的に共変
\item
  静止時のエネルギー \(E_0 = 8m\)（自然単位系で \(E_0 = 8\)）
\end{itemize}

\subsubsection{\texorpdfstring{2.2 反キンク解（巻き数
\(-1\)）}{2.2 反キンク解（巻き数 -1）}}\label{ux53cdux30adux30f3ux30afux89e3ux5dfbux304dux6570--1}

反キンク解は、\(\phi\) が \(2\pi\) から \(0\) へ遷移する局在波である
{[}2{]}。

\[\phi_{\bar{K}}(x, t) = 4\arctan\left(\exp\left(-\frac{x - vt}{\sqrt{1 - v^2}}\right)\right)\]

\textbf{性質:}

\begin{itemize}
\tightlist
\item
  \(x \to -\infty\) で \(\phi \to 2\pi\)、\(x \to +\infty\) で
  \(\phi \to 0\)
\item
  位相が \(2\pi\) だけ減少する ── \textbf{巻き数} \(N = -1\)
\item
  キンクと同じ質量・速度特性を持つ
\end{itemize}

\subsubsection{2.3
巻き数の定義と保存}\label{ux5dfbux304dux6570ux306eux5b9aux7fa9ux3068ux4fddux5b58}

\textbf{定義 2.1.} 場の配位 \(\phi(x, t)\)
の\textbf{巻き数}（トポロジカル荷）を以下で定義する {[}4, 10{]}。

\[N = \frac{1}{2\pi}\int_{-\infty}^{+\infty} \frac{\partial \phi}{\partial x}\,dx = \frac{\phi(+\infty) - \phi(-\infty)}{2\pi}\]

\textbf{命題 2.1（巻き数の保存）.}
sine-Gordon方程式の時間発展において、巻き数 \(N\) は保存される。

\emph{証明.} \(\phi\) の境界値 \(\phi(\pm\infty)\)
は真空値（\(2n\pi\)）に固定されており、連続的な時間発展によって \(2\pi\)
の整数倍だけ変化することはできない。したがって \(N\)
は時間に依存しない整数値として保存される。\(\square\)

この保存則は、ソリトンの「粒子性」------生成・消滅は対（キンク＋反キンク）でのみ可能------を保証する。

\subsubsection{2.4
ブリーザー解（束縛振動状態）}\label{ux30d6ux30eaux30fcux30b6ux30fcux89e3ux675fux7e1bux632fux52d5ux72b6ux614b}

ブリーザーはキンクと反キンクが束縛された振動解であり、巻き数 \(N = 0\)
を持つ {[}2, 11{]}。

\[\phi_B(x, t) = 4\arctan\left(\frac{\sqrt{1 - \omega^2}\;\cos(\omega t)}{\omega\;\cosh\left(\sqrt{1 - \omega^2}\;x\right)}\right)\]

ここで \(\omega\)（\(0 < \omega < 1\)）はブリーザーの内部振動数である。

\textbf{性質:}

\begin{itemize}
\tightlist
\item
  局在した振動解であり、空間的に指数関数的に減衰する
\item
  内部振動数 \(\omega\) が小さいほどブリーザーの空間的広がりが大きい
\item
  ブリーザーの静止質量は \(M_B = 16\sqrt{1 - \omega^2}\)
\item
  \(\omega \to 0\)
  の極限でブリーザーの質量はキンク-反キンク対の質量（\(2 \times 8 = 16\)）に近づく
\item
  \(\omega \to 1\) の極限でブリーザーの質量は \(0\)
  に近づき、線形の小振動（メソン）に移行する
\end{itemize}

\subsubsection{2.5
解の分類のまとめ}\label{ux89e3ux306eux5206ux985eux306eux307eux3068ux3081}

{\def\LTcaptype{none} % do not increment counter
\begin{longtable}[]{@{}ccccc@{}}
\toprule\noalign{}
解の種類 & 巻き数 \(N\) & 局在性 & 振動 & 質量 \\
\midrule\noalign{}
\endhead
\bottomrule\noalign{}
\endlastfoot
キンク & \(+1\) & 局在 & なし & \(8\) \\
反キンク & \(-1\) & 局在 & なし & \(8\) \\
ブリーザー & \(0\) & 局在 & あり & \(16\sqrt{1-\omega^2}\) \\
メソン（線形波） & \(0\) & 非局在 & あり & \(\geq 2\)（閾値） \\
\end{longtable}
}

\begin{center}\rule{0.5\linewidth}{0.5pt}\end{center}

\subsection{3.
ソリトン衝突と位相シフト}\label{ux30bdux30eaux30c8ux30f3ux885dux7a81ux3068ux4f4dux76f8ux30b7ux30d5ux30c8}

\subsubsection{3.1
キンク-キンク衝突}\label{ux30adux30f3ux30af-ux30adux30f3ux30afux885dux7a81}

同符号のキンク同士（\(N = +1\) と
\(N = +1\)）の衝突は反発的であり、2体解は以下で与えられる {[}2, 12{]}。

\[\phi_{KK}(x, t) = 4\arctan\left(\frac{v\,\sinh\left(\frac{x}{\sqrt{1-v^2}}\right)}{\cosh\left(\frac{vt}{\sqrt{1-v^2}}\right)}\right)\]

衝突後、2つのキンクは再び分離し、形状を完全に保存する。ただし、衝突がなかった場合と比較して、各キンクの位置が
\textbf{位相シフト} \(\Delta x\) だけずれる。

\[\Delta x = \frac{\sqrt{1 - v^2}}{v}\ln(v^2)\]

\subsubsection{3.2
キンク-反キンク衝突}\label{ux30adux30f3ux30af-ux53cdux30adux30f3ux30afux885dux7a81}

逆符号のキンクと反キンク（\(N = +1\) と
\(N = -1\)）の衝突は引力的であり、2体解は以下で与えられる {[}2, 9,
12{]}。

\[\phi_{K\bar{K}}(x, t) = 4\arctan\left(\frac{v^{-1}\,\cosh\left(\frac{vt}{\sqrt{1-v^2}}\right)}{\sinh\left(\frac{x}{\sqrt{1-v^2}}\right)}\right)\]

\textbf{性質:}

\begin{itemize}
\tightlist
\item
  十分なエネルギーがあれば、キンクと反キンクはすり抜けて再び分離する（完全弾性衝突：速度・形状は衝突前と同一、位相シフトのみ発生）
\item
  エネルギーが不十分な場合、ブリーザー（束縛状態）を形成する（永久に振動し続け、減衰しない）
\end{itemize}

\textbf{注記：} 古典的な完全可積分系としての
sine-Gordon方程式では、キンク-反キンク対の消滅→放射波（メソン）への変換は発生しない。これは量子効果または可積分性が破れた系（摂動・離散化・高次元化）でのみ生じる現象である。

\subsubsection{3.3
位相シフトの物理的意味}\label{ux4f4dux76f8ux30b7ux30d5ux30c8ux306eux7269ux7406ux7684ux610fux5473}

\textbf{命題 3.1（形状保存と位相シフト）.}
sine-Gordon方程式のソリトン衝突において、

\begin{enumerate}
\def\labelenumi{\arabic{enumi}.}
\tightlist
\item
  衝突後の各ソリトンの\textbf{形状}は衝突前と同一である。
\item
  衝突後の各ソリトンの\textbf{位置}は、衝突がなかった場合と比較して位相シフト
  \(\Delta x\) だけ変位する。
\item
  巻き数の総和は衝突前後で保存される。
\end{enumerate}

\emph{証明.} sine-Gordon方程式は逆散乱法によって厳密に解くことができ
{[}12{]}、多ソリトン解は漸近的に個々のソリトンの重ね合わせに分離する。各ソリトンの漸近的形状は衝突前と同一であるが、位相シフトが付加される。巻き数は位相的不変量であるため保存される。\(\square\)

この性質は本理論にとって決定的に重要である。ソリトン（＝粒子）の\textbf{種類}（巻き数＝配向構造）は相互作用で不変であり、\textbf{位置}（位相）のみが変化する。

\begin{center}\rule{0.5\linewidth}{0.5pt}\end{center}

\subsection{4.
完全可積分性と逆散乱法}\label{ux5b8cux5168ux53efux7a4dux5206ux6027ux3068ux9006ux6563ux4e71ux6cd5}

\subsubsection{4.1 ラックス対}\label{ux30e9ux30c3ux30afux30b9ux5bfe}

sine-Gordon方程式は\textbf{完全可積分系}であり、ラックス対（Lax
pair）を持つ {[}12, 13{]}。

光円錐座標 \(u = \frac{1}{2}(x - t)\), \(v = \frac{1}{2}(x + t)\)
を用いると、方程式は

\[\frac{\partial^2 \phi}{\partial u\,\partial v} = \sin\phi\]

と書ける。この方程式は以下のラックス対の両立条件として表される。

\[\Psi_u = U\Psi, \quad \Psi_v = V\Psi\]

ここで \(U\) と \(V\) はスペクトルパラメータ \(\lambda\) に依存する
\(2 \times 2\) 行列である {[}12{]}。

\subsubsection{4.2
無限個の保存量}\label{ux7121ux9650ux500bux306eux4fddux5b58ux91cf}

完全可積分性の帰結として、sine-Gordon方程式は無限個の独立な保存量を持つ
{[}12, 14{]}。最初の3つは以下の通りである。

\begin{itemize}
\item
  \textbf{エネルギー:}
  \[E = \int_{-\infty}^{+\infty}\left[\frac{1}{2}\left(\frac{\partial\phi}{\partial t}\right)^2 + \frac{1}{2}\left(\frac{\partial\phi}{\partial x}\right)^2 + (1-\cos\phi)\right]dx\]
\item
  \textbf{運動量:}
  \[P = -\int_{-\infty}^{+\infty}\frac{\partial\phi}{\partial t}\frac{\partial\phi}{\partial x}\,dx\]
\item
  \textbf{巻き数:}
  \[N = \frac{1}{2\pi}\int_{-\infty}^{+\infty}\frac{\partial\phi}{\partial x}\,dx\]
\end{itemize}

\subsubsection{4.3 Bäcklund変換}\label{buxe4cklundux5909ux63db}

sine-Gordon方程式は Bäcklund変換を持ち、既知の解から新しい解を生成できる
{[}7, 15{]}。

\(\phi_0\) を既知の解とするとき、以下の連立方程式の解 \(\phi_1\) もまた
sine-Gordon方程式の解である。

\[\frac{\partial}{\partial u}\left(\frac{\phi_1 + \phi_0}{2}\right) = a\,\sin\left(\frac{\phi_1 - \phi_0}{2}\right)\]

\[\frac{\partial}{\partial v}\left(\frac{\phi_1 - \phi_0}{2}\right) = \frac{1}{a}\sin\left(\frac{\phi_1 + \phi_0}{2}\right)\]

ここで \(a\) は Bäcklund変換のパラメータである。真空解 \(\phi_0 = 0\)
から出発して、1回の
Bäcklund変換でキンク解が、2回の変換でキンク-キンク解やブリーザー解が得られる
{[}15{]}。

\begin{center}\rule{0.5\linewidth}{0.5pt}\end{center}

\subsection{5.
高次元への拡張}\label{ux9ad8ux6b21ux5143ux3078ux306eux62e1ux5f35}

\subsubsection{5.1
2+1次元および3+1次元}\label{ux6b21ux5143ux304aux3088ux307331ux6b21ux5143}

sine-Gordon方程式の高次元への直接拡張は

\[\frac{\partial^2\phi}{\partial t^2} - \nabla^2\phi + \sin(\phi) = 0\]

であるが、2+1次元以上ではこの方程式は完全可積分ではなくなる
{[}16{]}。また、Derrickの定理 {[}17{]}
により、3+1次元ではスカラー場のみからなる安定な局在解（ソリトン）は存在しない。

\subsubsection{5.2
安定化の方法}\label{ux5b89ux5b9aux5316ux306eux65b9ux6cd5}

高次元で安定なソリトンを得るためには、以下の方法が知られている {[}10,
18{]}。

\begin{enumerate}
\def\labelenumi{\arabic{enumi}.}
\tightlist
\item
  \textbf{Skyrmeモデル}: 高次の微分項（Skyrme項）を追加する {[}8,
  18{]}。
\item
  \textbf{ゲージ場との結合}: スカラー場にゲージ場を結合させる。
\item
  \textbf{位相的安定化}:
  場が非自明なトポロジカル荷を持つ場合、エネルギー最小化がソリトンの崩壊を禁止する
  {[}10{]}。
\end{enumerate}

\subsubsection{5.3
閉空間上での安定化}\label{ux9589ux7a7aux9593ux4e0aux3067ux306eux5b89ux5b9aux5316}

高次元空間上でソリトンを扱う場合、Derrickの定理を回避する機構が必要である。周長
\(L\) の閉曲線 \(S^1\) 上では
Derrickの定理の前提（無限空間での漸近条件）が破れるため、安定なソリトン解が存在しうる
{[}10{]}。

\begin{center}\rule{0.5\linewidth}{0.5pt}\end{center}

\subsection{6. 物理的応用}\label{ux7269ux7406ux7684ux5fdcux7528}

\subsubsection{6.1
凝縮系物理学}\label{ux51ddux7e2eux7cfbux7269ux7406ux5b66}

\begin{itemize}
\tightlist
\item
  \textbf{ジョセフソン接合}:
  超伝導体間のトンネル接合における磁束量子の運動が
  sine-Gordon方程式で記述される
  {[}19{]}。キンクは磁束量子（フラクソン）に対応する。
\item
  \textbf{結晶転位}:
  Frenkel-Kontorova模型における結晶格子中の転位がキンク解に対応する
  {[}1{]}。
\end{itemize}

\subsubsection{6.2
素粒子物理学}\label{ux7d20ux7c92ux5b50ux7269ux7406ux5b66}

\begin{itemize}
\tightlist
\item
  \textbf{Skyrmeモデル}:
  バリオンが位相的ソリトン（スカーミオン）として記述される {[}8,
  18{]}。バリオン数はトポロジカル荷に対応する。
\item
  \textbf{量子sine-Gordon模型と大規模Thirring模型の等価性}: Coleman
  {[}20{]} および Mandelstam {[}21{]}
  は、sine-Gordon模型の量子論が大規模
  Thirring模型（フェルミオンの場の理論）と等価であることを示した。ボソン的ソリトンがフェルミオンに対応するという
  \textbf{ボソン-フェルミオン対応}
  は、1+1次元場の理論の最も顕著な結果の一つである。
\end{itemize}

\subsubsection{6.3 微分幾何学}\label{ux5faeux5206ux5e7eux4f55ux5b66}

\begin{itemize}
\tightlist
\item
  \textbf{擬球面}:
  定負曲率（\(K = -1\)）の曲面の第一基本形式と第二基本形式の両立条件が
  sine-Gordon方程式に帰着される {[}5,
  6{]}。キンク解は擬球面上の測地線に対応する。
\end{itemize}

\begin{center}\rule{0.5\linewidth}{0.5pt}\end{center}

\subsection{7.
物理的解釈への示唆}\label{ux7269ux7406ux7684ux89e3ux91c8ux3078ux306eux793aux5506}

\subsubsection{7.1
ソリトンの粒子的性質のまとめ}\label{ux30bdux30eaux30c8ux30f3ux306eux7c92ux5b50ux7684ux6027ux8ceaux306eux307eux3068ux3081}

sine-Gordon方程式のソリトン解は、以下の粒子的性質を自然に備えている。

{\def\LTcaptype{none} % do not increment counter
\begin{longtable}[]{@{}cc@{}}
\toprule\noalign{}
ソリトンの性質 & 粒子との対応 \\
\midrule\noalign{}
\endhead
\bottomrule\noalign{}
\endlastfoot
巻き数 \(N = \pm 1\) & 粒子／反粒子 \\
巻き数の保存 & 粒子数の保存 \\
衝突時の形状保存 & 粒子の同一性 \\
位相シフト（衝突後の位置変化） & 相互作用による散乱 \\
ローレンツ共変性 & 特殊相対論的構造 \\
ブリーザー（\(N = 0\)の束縛状態） & メソン的束縛状態 \\
ポテンシャルの周期性（\(2\pi\)） & 位相の離散性 \\
ボソン-フェルミオン対応 {[}20, 21{]} & スピン-統計 \\
\end{longtable}
}

\subsubsection{7.2
矩形波記述の意義}\label{ux77e9ux5f62ux6ce2ux8a18ux8ff0ux306eux610fux7fa9}

§8--§9 で示すように、これらの粒子的性質の全ては矩形波極限
\(\sigma \to 0\)
において保持される。このことは、ソリトンの本質的な情報が連続場の滑らかな形状にあるのではなく、位相的・散乱的な離散構造に集約されていることを意味する。

\begin{center}\rule{0.5\linewidth}{0.5pt}\end{center}

\subsection{8.
ソリトンの有限パラメータ表現}\label{ux30bdux30eaux30c8ux30f3ux306eux6709ux9650ux30d1ux30e9ux30e1ux30fcux30bfux8868ux73fe}

\subsubsection{8.1
動機：連続場から有限個のパラメータへ}\label{ux52d5ux6a5fux9023ux7d9aux5834ux304bux3089ux6709ux9650ux500bux306eux30d1ux30e9ux30e1ux30fcux30bfux3078}

sine-Gordon方程式の解は連続場 \(\phi(x, t)\)
として定義されるが、ソリトン解はその局在性ゆえに、\textbf{有限個のパラメータで完全に特徴づけられる}。これは偶然ではない。ソリトンは非線形場のモード分解における\textbf{有限モード解}であり、無限自由度の場の理論の中で有限次元の不変多様体（moduli
space）上に制限された解である {[}10, 22{]}。

本節では、sine-Gordon方程式のソリトン解を有限個のパラメータで記述する定式化を与え、連続的な波動方程式の解がいかに少数の離散的な情報に縮約されるかを明示する。

\subsubsection{8.2
キンクの5パラメータ表現}\label{ux30adux30f3ux30afux306e5ux30d1ux30e9ux30e1ux30fcux30bfux8868ux73fe}

§2.1 のキンク解を再掲する。

\[\phi_K(x, t) = 4\arctan\left(\exp\left(\frac{x - x_0 - vt}{\ell}\right)\right)\]

ここで \(\ell = \sqrt{1 - v^2}\)
はキンクの幅パラメータである。この解は以下の5つのパラメータで\textbf{完全に}決定される。

\textbf{定義 8.1（ソリトンパラメータ）.} 1+1次元 sine-Gordon
キンクを特徴づける5つのパラメータを以下のように定義する。

{\def\LTcaptype{none} % do not increment counter
\begin{longtable}[]{@{}
  >{\centering\arraybackslash}p{(\linewidth - 6\tabcolsep) * \real{0.2500}}
  >{\centering\arraybackslash}p{(\linewidth - 6\tabcolsep) * \real{0.2500}}
  >{\centering\arraybackslash}p{(\linewidth - 6\tabcolsep) * \real{0.2500}}
  >{\centering\arraybackslash}p{(\linewidth - 6\tabcolsep) * \real{0.2500}}@{}}
\toprule\noalign{}
\begin{minipage}[b]{\linewidth}\centering
記号
\end{minipage} & \begin{minipage}[b]{\linewidth}\centering
名称
\end{minipage} & \begin{minipage}[b]{\linewidth}\centering
定義
\end{minipage} & \begin{minipage}[b]{\linewidth}\centering
値域
\end{minipage} \\
\midrule\noalign{}
\endhead
\bottomrule\noalign{}
\endlastfoot
\(N\) & 巻き数 & トポロジカル荷 \(\pm 1\) & \(\{-1, +1\}\) \\
\(A\) & 最大波高 & \(\phi\) の全変化量 \(= 2\pi |N|\) &
\(2\pi\)（固定） \\
\(\theta_0\) & 中心位相 & ソリトン中心の空間位相 \(= 2\pi x_0 / L\) &
\([0, 2\pi)\) \\
\(\alpha\) & 移動位相 & 速度の位相角表現 \(v = \sin\alpha\) &
\((-\pi/2, \pi/2)\) \\
\(\sigma\) & 広がり位相 & 幅の位相角表現 \(\ell = \cos\alpha\) &
\((0, 1]\) \\
\end{longtable}
}

\textbf{命題 8.1（パラメータの完全性）.}
sine-Gordon方程式の単一キンク（反キンク）解は、上記5パラメータ
\((N, A, \theta_0, \alpha, \sigma)\)
により一意に決定される。逆に、任意のキンク解から上記5パラメータが一意に読み取れる。

\emph{証明.} キンク解の一般形は
\(\phi = 4\arctan(\exp(\pm(x - x_0 - vt)/\sqrt{1-v^2}))\) であり、符号
\(\pm\) が \(N\)、変位 \(x_0\) が \(\theta_0\)、速度 \(v\) が \(\alpha\)
を決定する。\(A = 2\pi\) は巻き数 \(|N| = 1\)
から固定され、\(\sigma = \cos\alpha\) は \(\alpha\)
から従属的に決まる。自由パラメータは \((N, \theta_0, \alpha)\)
の3つであり、\(A\) と \(\sigma\) はこれらから導出される。\(\square\)

\textbf{注記.}
自由パラメータは実質3つ（\(N, \theta_0, \alpha\)）であるが、物理的解釈の明確さのために5パラメータ表現を採用する。\(A\)
は巻き数から、\(\sigma\) は移動位相から決まる従属パラメータである。

\subsubsection{8.3
移動位相と広がり位相の関係}\label{ux79fbux52d5ux4f4dux76f8ux3068ux5e83ux304cux308aux4f4dux76f8ux306eux95a2ux4fc2}

移動位相 \(\alpha\) と広がり位相 \(\sigma\)
の間には以下の恒等式が成り立つ。

\[\sin^2\alpha + \sigma^2 = v^2 + \ell^2 = 1\]

すなわち、\((\sin\alpha, \sigma) = (v, \ell)\) は単位円上の点である。

\textbf{物理的解釈:}

\begin{itemize}
\tightlist
\item
  \(\alpha = 0\)（静止）のとき \(\sigma = 1\)：ソリトンは最大幅を持つ
\item
  \(\alpha \to \pm\pi/2\)（光速）のとき
  \(\sigma \to 0\)：ソリトンの幅が0に収縮する（ローレンツ収縮）
\item
  移動が速いほど空間的に狭くなる ── \textbf{運動と形状が一つの位相角
  \(\alpha\) で統一的に記述される}
\end{itemize}

\subsubsection{8.4
ブリーザーの6パラメータ表現}\label{ux30d6ux30eaux30fcux30b6ux30fcux306e6ux30d1ux30e9ux30e1ux30fcux30bfux8868ux73fe}

§2.4 のブリーザー解は、キンクの5パラメータに加えて内部振動数 \(\omega\)
を必要とする。

\textbf{定義 8.2（ブリーザーパラメータ）.}

{\def\LTcaptype{none} % do not increment counter
\begin{longtable}[]{@{}
  >{\centering\arraybackslash}p{(\linewidth - 6\tabcolsep) * \real{0.2500}}
  >{\centering\arraybackslash}p{(\linewidth - 6\tabcolsep) * \real{0.2500}}
  >{\centering\arraybackslash}p{(\linewidth - 6\tabcolsep) * \real{0.2500}}
  >{\centering\arraybackslash}p{(\linewidth - 6\tabcolsep) * \real{0.2500}}@{}}
\toprule\noalign{}
\begin{minipage}[b]{\linewidth}\centering
記号
\end{minipage} & \begin{minipage}[b]{\linewidth}\centering
名称
\end{minipage} & \begin{minipage}[b]{\linewidth}\centering
定義
\end{minipage} & \begin{minipage}[b]{\linewidth}\centering
値域
\end{minipage} \\
\midrule\noalign{}
\endhead
\bottomrule\noalign{}
\endlastfoot
\(N\) & 巻き数 & \(0\)（キンク-反キンク束縛） & \(\{0\}\)（固定） \\
\(A\) & 最大波高 & \(4\arctan(\sqrt{1-\omega^2}/\omega)\) &
\((0, \pi)\) \\
\(\theta_0\) & 中心位相 & ブリーザー中心の空間位相 & \([0, 2\pi)\) \\
\(\alpha\) & 移動位相 & 全体の運動 \(v = \sin\alpha\) &
\((-\pi/2, \pi/2)\) \\
\(\sigma\) & 広がり位相 & \(1/\sqrt{1-\omega^2}\) & \([1, \infty)\) \\
\(\omega\) & 内部位相速度 & 内部振動の角振動数 & \((0, 1)\) \\
\end{longtable}
}

ブリーザーは \(N = 0\) であるが、内部振動 \(\omega\)
という追加の自由度を持つ。

\subsubsection{8.5
多ソリトン状態のパラメータ表現}\label{ux591aux30bdux30eaux30c8ux30f3ux72b6ux614bux306eux30d1ux30e9ux30e1ux30fcux30bfux8868ux73fe}

\(n\)
個のソリトンからなる多体状態は、各ソリトンのパラメータの集合として表現される
{[}10, 12{]}。

\textbf{定義 8.3（\(n\)-ソリトン状態）.} \(n\)
個のソリトンからなる状態を以下のパラメータ集合で定義する。

\[\mathcal{S}_n = \{(N_i, A_i, \theta_{0,i}, \alpha_i, \sigma_i)\}_{i=1}^{n}\]

\textbf{命題 8.2（漸近的完全性）.} 十分に離れた \(n\)
個のソリトンからなる場の配位 \(\phi(x, t)\) は、パラメータ集合
\(\mathcal{S}_n\)
により漸近的に完全に決定される。ソリトン間の重なりが無視できる領域では、

\[\phi(x, t) \approx \sum_{i=1}^{n} \phi_i(x, t; N_i, \theta_{0,i}, \alpha_i)\]

が成り立つ。ここで \(\phi_i\) は各ソリトンの単体解である。

\emph{証明.} 逆散乱法 {[}12{]} により、\(n\)-ソリトン解は \(n\)
個の散乱データ（固有値とノルミング定数）で完全に決定される。漸近領域では散乱データと上記パラメータは1対1に対応する。\(\square\)

\subsubsection{8.6
衝突のパラメータ空間での記述}\label{ux885dux7a81ux306eux30d1ux30e9ux30e1ux30fcux30bfux7a7aux9593ux3067ux306eux8a18ux8ff0}

§3 のソリトン衝突は、パラメータ空間上の写像として記述できる。

\textbf{命題 8.3（衝突写像）.} 2つのソリトン
\(\mathcal{S}_1 = (N_1, \theta_{0,1}, \alpha_1)\) と
\(\mathcal{S}_2 = (N_2, \theta_{0,2}, \alpha_2)\)
の衝突は、以下のパラメータ変換として記述される。

\[\mathcal{S}_1' = (N_1,\; \theta_{0,1} + \Delta\theta_1,\; \alpha_1)\]
\[\mathcal{S}_2' = (N_2,\; \theta_{0,2} + \Delta\theta_2,\; \alpha_2)\]

すなわち：

\begin{enumerate}
\def\labelenumi{\arabic{enumi}.}
\tightlist
\item
  \textbf{巻き数 \(N_i\) は不変} ── 粒子の種類は変化しない
\item
  \textbf{移動位相 \(\alpha_i\) は不変} ── 速度（運動量）は保存される
\item
  \textbf{中心位相 \(\theta_{0,i}\) のみが変化} ── 位相シフト
  \(\Delta\theta_i\) だけ変位する
\end{enumerate}

\textbf{系 8.1.} ソリトン衝突の全情報は、位相シフト \(\Delta\theta_i\)
に集約される。位相シフトは入射パラメータから計算可能であり {[}12{]}、

\[\Delta\theta_1 = -\Delta\theta_2 = f(N_1, N_2, \alpha_1 - \alpha_2)\]

の形を取る。すなわち、衝突は相対移動位相 \(\alpha_1 - \alpha_2\)
のみに依存する。

\subsubsection{8.7
連続場からパラメータ表現への縮約}\label{ux9023ux7d9aux5834ux304bux3089ux30d1ux30e9ux30e1ux30fcux30bfux8868ux73feux3078ux306eux7e2eux7d04}

以上をまとめると、sine-Gordon方程式の解の空間は以下のように階層化される。

{\def\LTcaptype{none} % do not increment counter
\begin{longtable}[]{@{}
  >{\centering\arraybackslash}p{(\linewidth - 4\tabcolsep) * \real{0.4286}}
  >{\centering\arraybackslash}p{(\linewidth - 4\tabcolsep) * \real{0.2857}}
  >{\centering\arraybackslash}p{(\linewidth - 4\tabcolsep) * \real{0.2857}}@{}}
\toprule\noalign{}
\begin{minipage}[b]{\linewidth}\centering
記述レベル
\end{minipage} & \begin{minipage}[b]{\linewidth}\centering
自由度
\end{minipage} & \begin{minipage}[b]{\linewidth}\centering
情報量
\end{minipage} \\
\midrule\noalign{}
\endhead
\bottomrule\noalign{}
\endlastfoot
連続場 \(\phi(x, t)\) & 無限次元 & 各時空点で実数値 \\
\(n\)-ソリトン解 & \(3n\) 次元 & 各ソリトン
\((N_i, \theta_{0,i}, \alpha_i)\) \\
衝突後の変化 & \(n\) 次元 & 位相シフト \(\Delta\theta_i\) のみ \\
\end{longtable}
}

\textbf{命題 8.4（ソリトンセクターの有限次元性）.}
sine-Gordon方程式の解の空間のうち、\(n\)
個のソリトンを含むセクターは、\(3n\)
次元の有限次元多様体（モジュライ空間）\(\mathcal{M}_n\) として記述される
{[}10,
22{]}。連続場の無限自由度は、ソリトンセクターにおいて有限個のパラメータに\textbf{厳密に}縮約される。

\emph{証明.} 逆散乱法において、\(n\)-ソリトンセクターは \(n\)
個の離散固有値 \(\zeta_i\) とそれに付随するノルミング定数 \(c_i\)
で特徴づけられる {[}12, 13{]}。各固有値は巻き数 \(N_i\) と速度
\(v_i\)（= \(\sin\alpha_i\)）を、ノルミング定数は初期位置 \(x_{0,i}\)（=
\(\theta_{0,i}\)）を決定する。放射成分（連続スペクトル）を除けば、ソリトンセクターの自由度は正確に
\(3n\) である。\(\square\)

この縮約は近似ではない。ソリトンセクターにおいて、無限次元の偏微分方程式は\textbf{厳密に}有限次元の常微分方程式系に帰着される。

\subsubsection{8.8
矩形波との関係}\label{ux77e9ux5f62ux6ce2ux3068ux306eux95a2ux4fc2}

キンク解の形状を改めて観察する。

\[\phi_K(x) = 4\arctan(e^{x/\ell}) = 2\pi \cdot \frac{1}{1 + e^{-x/\ell}}\]

右辺はロジスティック関数（シグモイド関数）であり、\(\ell \to 0\)
の極限で\textbf{階段関数}（矩形波の半周期）に収束する。

\[\lim_{\ell \to 0} \phi_K(x) = \begin{cases} 0 & (x < 0) \\ 2\pi & (x > 0) \end{cases}\]

\textbf{命題 8.5（矩形波極限）.} キンクの広がりパラメータ
\(\sigma = \ell \to 0\)（すなわち
\(\alpha \to \pm\pi/2\)、光速極限）において、キンク解は矩形波（階段関数）に収束する。逆に、有限幅
\(\sigma > 0\) のキンクは、矩形波の\textbf{平滑化}として理解できる。

\emph{証明.} \(\phi_K(x) = 2\pi/(1 + e^{-x/\ell})\)
はロジスティック関数であり、\(\ell \to 0\) で Heaviside階段関数
\(2\pi \cdot H(x)\) に各点収束する。\(\square\)

\textbf{系 8.2.}
ソリトンは矩形波の有限モード近似として解釈できる。連続的な \(\tanh\)
型の形状は、離散的な階段関数の平滑化バージョンに過ぎない。広がりパラメータ
\(\sigma\) は平滑化の程度を制御する。

この観点から、ソリトンのパラメータ表現は以下のように再解釈される。

{\def\LTcaptype{none} % do not increment counter
\begin{longtable}[]{@{}ccc@{}}
\toprule\noalign{}
パラメータ & 連続場での意味 & 矩形波（離散）での意味 \\
\midrule\noalign{}
\endhead
\bottomrule\noalign{}
\endlastfoot
\(N = \pm 1\) & 巻き数 & 階段の向き（上昇/下降） \\
\(A = 2\pi\) & 位相の全変化量 & 階段の高さ \\
\(\theta_0\) & ソリトン中心の位置 & 階段のエッジ位置 \\
\(\alpha\) & 速度の位相角 & エッジの移動方向 \\
\(\sigma\) & 平滑化の幅 & \textbf{\(\sigma \to 0\) で矩形波に一致} \\
\end{longtable}
}

\begin{center}\rule{0.5\linewidth}{0.5pt}\end{center}

\subsection{9.
矩形波極限における諸性質の保存}\label{ux77e9ux5f62ux6ce2ux6975ux9650ux306bux304aux3051ux308bux8af8ux6027ux8ceaux306eux4fddux5b58}

\subsubsection{9.1 目的}\label{ux76eeux7684}

§1〜§7 で述べた sine-Gordon方程式の諸性質は、連続場 \(\phi(x, t)\)
に対して証明されている。§8.8
で示したように、ソリトン解は広がりパラメータ \(\sigma \to 0\)
の極限で矩形波（階段関数）に収束する。

本節では、§1〜§7 の各性質について、(a)
ソリトン解に対して成立すること（これは厳密解であるから自明）、(b)
矩形波極限 \(\sigma \to 0\) においても保持されることを体系的に検証する。

もし全性質が矩形波極限で保持されるならば、矩形波記述は「都合のよい近似」ではなく、ソリトンの持つ構造を\textbf{完全に継承する厳密な極限}であると結論できる。

\subsubsection{9.2
検証の枠組み}\label{ux691cux8a3cux306eux67a0ux7d44ux307f}

各性質を以下の3段階で検証する。

\begin{enumerate}
\def\labelenumi{\arabic{enumi}.}
\tightlist
\item
  \textbf{ソリトン解での成立}：厳密解として sine-Gordon方程式を満たすか
\item
  \textbf{\(\sigma\)-依存性の分析}：その性質が広がりパラメータ
  \(\sigma\) に依存するか
\item
  \textbf{矩形波極限の存在}：\(\sigma \to 0\)
  で性質が保持されるか、発散するか
\end{enumerate}

\subsubsection{9.3 巻き数の定義と保存（§2.3 →
矩形波）}\label{ux5dfbux304dux6570ux306eux5b9aux7fa9ux3068ux4fddux5b582.3-ux77e9ux5f62ux6ce2}

\textbf{§2.3 の性質：} 巻き数
\(N = [\phi(+\infty) - \phi(-\infty)] / 2\pi\) は時間発展で保存される。

\textbf{ソリトン解での成立：} キンク解は \(\phi(-\infty) = 0\),
\(\phi(+\infty) = 2\pi\) を与え、\(N = +1\) である。これは
sine-Gordon方程式の厳密解であるから、\(N\) の保存は自明。\(\checkmark\)

\textbf{\(\sigma\)-依存性：} 巻き数の定義
\(N = [\phi(+\infty) - \phi(-\infty)] / 2\pi\)
は\textbf{境界値のみ}に依存し、\(\phi\) の内部形状（したがって
\(\sigma\)）に一切依存しない。

\textbf{命題 9.1（巻き数の \(\sigma\)-非依存性）.} 巻き数 \(N\)
は広がりパラメータ \(\sigma\) に依存しない。矩形波極限 \(\sigma \to 0\)
においても \(N\) は同一の整数値を保つ。

\emph{証明.} 任意の \(\sigma > 0\)
に対して、\(\phi_K^\sigma(\pm\infty)\) は真空値 \(0, 2\pi\)
に固定されている。\(\sigma \to 0\) の極限で
\(\phi_K^\sigma(x) \to 2\pi \cdot H(x)\)
に各点収束するが、\(H(-\infty) = 0\), \(H(+\infty) = 1\) であるから
\(N = (2\pi \cdot 1 - 2\pi \cdot 0)/2\pi = 1\) は不変。\(\square\)

\textbf{矩形波での巻き数：} 矩形波 \(\phi(x) = 2\pi H(x - x_0)\)
に対して \(N = +1\)。階段の向きが上昇なら \(N = +1\)、下降なら
\(N = -1\)。巻き数はまさに階段の符号そのものである。\(\checkmark\)

\subsubsection{9.4 巻き数の保存則（命題 2.1 →
矩形波）}\label{ux5dfbux304dux6570ux306eux4fddux5b58ux5247ux547dux984c-2.1-ux77e9ux5f62ux6ce2}

\textbf{§2.3 の性質：} \(N\)
は時間発展で保存される。生成・消滅はキンク-反キンク対でのみ可能。

\textbf{\(\sigma\)-依存性：}
保存則の証明（§2.3）は「境界値が真空値に固定され、連続的な時間発展で
\(2\pi\) の整数倍だけ変化できない」ことのみに依拠する。この論法は
\(\sigma\) に依存しない。

\textbf{矩形波での保存：} 矩形波の境界値は \(0\) と \(2\pi\)
であり、離散的な値のみを取る。離散値は連続変形では変化しない。したがって、矩形波における巻き数の保存は、連続場の場合よりも\textbf{むしろ自明}である。

\textbf{命題 9.2.}
矩形波極限において巻き数の保存則は厳密に成立する。さらに、矩形波では
\(\phi\) 自体が離散値（\(0\) または
\(2\pi\)）のみを取るため、保存は位相的のみならず\textbf{値の離散性}からも保証される。\(\checkmark\)

\subsubsection{9.5 キンク・反キンク解の存在（§2.1, §2.2 →
矩形波）}\label{ux30adux30f3ux30afux53cdux30adux30f3ux30afux89e3ux306eux5b58ux57282.1-2.2-ux77e9ux5f62ux6ce2}

\textbf{§2.1--2.2 の性質：} 巻き数 \(\pm 1\) の局在解が存在し、速度
\(v\) で伝播する。

\textbf{ソリトン解での成立：} キンク
\(\phi_K = 4\arctan(\exp((x-vt)/\sigma))\) は厳密解。\(\checkmark\)

\textbf{矩形波極限：} \(\sigma \to 0\) で
\(\phi_K \to 2\pi \cdot H(x - vt)\)。これは速度 \(v\)
で移動する階段関数であり、局在性と伝播性を保持する。

\textbf{命題 9.3（矩形波キンクの存在）.} 矩形波キンク
\(\phi(x, t) = 2\pi \cdot H(x - vt)\) は以下の性質を持つ。 1. 巻き数
\(N = +1\) 2. 速度 \(v\) で移動 3. 空間の1点（\(x = vt\)）に局在 4.
sine-Gordon方程式の\textbf{超関数解}（distributional
solution）として成立する

\emph{証明.} (1)--(3) は定義から明らか。(4)
について：\(\phi = 2\pi H(x - vt)\) の超関数微分は
\(\phi_x = 2\pi\delta(x-vt)\), \(\phi_{xx} = 2\pi\delta'(x-vt)\),
\(\phi_{tt} = 2\pi v^2\delta'(x-vt)\)
であり、\(\sin(\phi) = \sin(2\pi H(x-vt)) = 0\)（\(x \neq vt\) で
\(\phi = 0\) or \(2\pi\), いずれも \(\sin = 0\)）。波動方程式部分
\(\phi_{tt} - \phi_{xx} = 2\pi(v^2-1)\delta'(x-vt)\) は
\(v = \pm 1\)（光速）の場合にのみ消滅する。\(v \neq 1\)
の場合、矩形波は古典解ではなく、有限幅 \(\sigma > 0\)
の平滑化が必要となる。\(\square\)

\textbf{注記.} \(v \neq 1\) で矩形波が厳密解でないことは、§8.3 の関係
\(\sigma = \cos\alpha = \sqrt{1-v^2}\) と整合する。\(v < 1\) ならば
\(\sigma > 0\)
であり、ソリトンは有限の幅を持つ。矩形波（\(\sigma = 0\)）は
\(v = 1\)（光速）にのみ対応する。光速以外の粒子が有限の「広がり」を持つことは、量子力学における波束の広がりと定性的に一致する。

\subsubsection{9.6 ブリーザー解（§2.4 →
矩形波）}\label{ux30d6ux30eaux30fcux30b6ux30fcux89e32.4-ux77e9ux5f62ux6ce2}

\textbf{§2.4 の性質：}
キンク-反キンク束縛状態としてのブリーザーが存在する。

\textbf{矩形波極限：} ブリーザーは巻き数 \(N = 0\)
の振動解である。矩形波極限では、これは \(0\) と \(2\pi\)
の間を振動する\textbf{矩形パルス}に対応する。

\textbf{命題 9.4.}
ブリーザーの矩形波極限は、空間的に局在した矩形パルスの振動である。パルスの存在（\(\phi = 2\pi\)）と不在（\(\phi = 0\)）の2状態間を周期的に遷移する。質量
\(M_B = 16\sqrt{1-\omega^2}\) は \(\omega\) に依存し、\(\sigma\)
には直接依存しない。\(\checkmark\)

\subsubsection{9.7 ソリトン衝突と位相シフト（§3 →
矩形波）}\label{ux30bdux30eaux30c8ux30f3ux885dux7a81ux3068ux4f4dux76f8ux30b7ux30d5ux30c83-ux77e9ux5f62ux6ce2}

\textbf{§3 の性質：} 衝突後に (i) 形状保存、(ii) 位相シフト、(iii)
巻き数保存。

これは最も重要な性質であり、各項を個別に検証する。

\textbf{(i) 形状保存：}

\textbf{ソリトン解：} 衝突後のソリトンは衝突前と同一の形状（同じ
\(\sigma\)）を持つ {[}12{]}。\(\checkmark\)

\textbf{矩形波極限：} 矩形波は \(\sigma = 0\)
であり、形状パラメータは「階段の高さ \(2\pi\)」のみ。衝突前後で
\(\sigma = 0\)
が保たれるならば、形状保存は自動的に成立する。そして、ソリトン衝突において
\(\sigma\) は保存量である（命題 8.3：\(\alpha_i\) 不変 \(\Rightarrow\)
\(\sigma_i = \cos\alpha_i\) 不変）ため、\(\sigma = 0\)
も保存される。\(\checkmark\)

\textbf{(ii) 位相シフト：}

\textbf{ソリトン解：} 衝突後の位相シフトは
\(\Delta\theta = f(N_1, N_2, \alpha_1 - \alpha_2)\)
{[}12{]}。\(\checkmark\)

\textbf{矩形波極限：} 位相シフトの公式は
\(\alpha_i\)（移動位相）のみに依存し、\(\sigma_i\)（広がり）には依存しない（§8.6
系8.1参照）。したがって \(\sigma \to 0\)
でも位相シフトの値は変化しない。

\textbf{命題 9.5（位相シフトの \(\sigma\)-非依存性）.}
ソリトン衝突の位相シフト \(\Delta\theta\)
は、各ソリトンの広がりパラメータ \(\sigma_i\) に依存しない。矩形波極限
\(\sigma \to 0\)
においても、位相シフトはソリトンの場合と同一の値を取る。

\emph{証明.} 逆散乱法において、位相シフトは散乱データの離散固有値
\(\zeta_i\) のみから計算される {[}12, 13{]}。\(\zeta_i\)
はソリトンの速度 \(v_i = \sin\alpha_i\) と巻き数 \(N_i\)
で決定され、ノルミング定数（位置情報）や放射成分とは独立である。ソリトンの幅
\(\sigma_i\)
はノルミング定数に関係するが、位相シフトには寄与しない。\(\square\)

\textbf{(iii) 巻き数保存：} §9.4 で既に示した。\(\checkmark\)

\subsubsection{9.8 完全可積分性と保存量（§4 →
矩形波）}\label{ux5b8cux5168ux53efux7a4dux5206ux6027ux3068ux4fddux5b58ux91cf4-ux77e9ux5f62ux6ce2}

\textbf{§4 の性質：} エネルギー、運動量、巻き数の保存。

\textbf{エネルギー：}

ソリトンのエネルギーは \(E = 8/\sigma = 8/\cos\alpha = 8/\sqrt{1-v^2}\)
である。\(\sigma \to 0\) でエネルギーは発散する。

\textbf{命題 9.6（エネルギーの \(\sigma\)-依存性）.}
ソリトンのエネルギーは \(\sigma\) に依存し、\(\sigma \to 0\) で
\(E \to \infty\)
に発散する。これは矩形波極限が無限エネルギーを持つことを意味する。

\textbf{命題 9.6a（本理論の適用限界）.}
sine-Gordonソリトン理論は、以下の暗黙の前提に基づいている。

\begin{enumerate}
\def\labelenumi{\arabic{enumi}.}
\tightlist
\item
  \textbf{起点の固定}: 場の初期値は真空値 \(\phi(-\infty) = 0\)
  に固定される（§2.1）。
\item
  \textbf{ポテンシャルの周期性}: \(V(\phi) = 1-\cos\phi\) は周期
  \(2\pi\) を持ち、ターゲット空間は実質的に \(S^1\)
  である。ポテンシャルの「一周」の大きさ（\(2\pi\)）は理論に固定されており、変更の自由度を持たない。
\item
  \textbf{終端の真空値への強制}: 無限空間 \((-\infty, +\infty)\)
  上での有限エネルギーの要請により、\(\phi(+\infty) = 2n\pi\)（\(n \in \mathbb{Z}\)）が強制される。これは
  \(V(2n\pi) = 0\)（ポテンシャル最小）であるためであり、\(\phi = (2n+1)\pi\)
  は
  \(V = 2\)（ポテンシャル最大）であるため、無限空間上ではエネルギーが発散し排除される。
\item
  \textbf{端数巻き数の排除}: 上記の帰結として、巻き数 \(N\)
  は常に整数となる。半整数巻き数（\(\phi(+\infty) = (2n+1)\pi\)）は有限エネルギーの要請により自動的に排除される。
\item
  \textbf{境界反射の不在}:
  空間が無限または閉（\(S^1\)）であるため、物理的な境界が存在せず、ソリトンの境界反射は発生しない。
\item
  \textbf{完全弾性衝突と相互作用の不在}:
  本理論は完全可積分系（§4）であり、無限個の保存量を持つ。この過剰な拘束の帰結として、ソリトン衝突は完全弾性衝突として記述され、衝突前後でソリトンの速度・形状・エネルギーは\textbf{個別に}保存される（§3.3）。運動量の保存則を破らないために、変化するのは各ソリトンの位置（位相シフト
  \(\Delta x\)）のみであり、速度の交換・形状の変形・エネルギーの散逸は一切発生しない。これは事実上、ソリトン同士が\textbf{相互作用しない自由粒子}として振る舞うことを意味する。同一粒子（例：キンク-キンク衝突）の場合、「すり抜けたのか反発したのか」は原理的に区別不能である。
\end{enumerate}

\emph{帰結.}
この理論が整合的に適用されるのは、以下の2つの場合に限られる。

\begin{itemize}
\tightlist
\item
  \textbf{無限空間}
  \((-\infty, +\infty)\)：終端が無限遠にあり、到達しない。エネルギーは各キンクにつき
  \(E = 8\)（静止時）で有限。
\item
  \textbf{閉曲線} \(S^1\)（周長 \(L\)）：周期的境界条件
  \(\phi(x+L) = \phi(x) + 2\pi N\)。境界が存在しない。
\end{itemize}

\textbf{開いた弦}（有限区間 \([0, L]\) 上の Dirichlet / Neumann
境界条件）は本理論の適用範囲外である。その理由は、有限区間 \([0, L]\)
は可縮空間（\(\pi_1([0,L]) = 0\)）であり、巻き数が位相的不変量として定義されず、ソリトンの位相的安定性が保証されないためである。さらに有限区間上では
\(\phi(L) = (2n+1)\pi\)（半整数巻き数、ポテンシャル最大値）がエネルギー的に許容されるため、整数巻き数への制限が成立しない。

\textbf{エネルギーの自己相殺構造：} 各 \(2\pi\)
回転（1キンク）のエネルギーは
\(E = 8\)（静止時）で自己完結する。これはポテンシャル
\(V(\phi) = 1-\cos\phi\) が一周期で \(0 \to 0\)
に戻るためである。キンク-反キンク対（\(4\pi\) 回転、巻き数
\(N = 0\)）は対消滅によりエネルギーをゼロに戻すことができる。観測可能な永続的エネルギーは、対消滅できない余りのキンク（\(|N|\)
個）の寄与 \(8|N|\) のみである。

\textbf{運動量と巻き数：} 運動量 \(P = -Ev\)
も同様の構造を持つ。巻き数は §9.3 で示した通り
\(\sigma\)-非依存。\(\checkmark\)

\subsubsection{9.9 高次元への拡張（§5 →
矩形波）}\label{ux9ad8ux6b21ux5143ux3078ux306eux62e1ux5f355-ux77e9ux5f62ux6ce2}

\textbf{§5 の性質：}
Derrickの定理により、3+1次元の無限空間では安定な局在解が存在しない。

\textbf{矩形波極限への影響：}
Derrickの定理の前提は\textbf{無限空間における有限エネルギー解}の存在である
{[}17{]}。閉曲線 \(S^1\)（周長
\(L\)）上では、以下の2つの理由で定理の前提が破れる。

\begin{enumerate}
\def\labelenumi{\arabic{enumi}.}
\tightlist
\item
  \textbf{空間が閉じている}: スケール変換
  \(\phi(x) \to \phi(\lambda x)\) における \(\lambda\)
  の範囲が制限される
\item
  \textbf{位相的保護}: 巻き数 \(N \in \pi_1(S^1) = \mathbb{Z}\)
  が位相的不変量として保護される
\end{enumerate}

\textbf{命題 9.7.} 閉曲線 \(S^1\)（周長
\(L\)）上では、矩形波（\(\sigma = 0\)）を含むあらゆる幅の局在解が位相的に安定である。Derrickの定理は閉空間には適用されない。\(\checkmark\)

\subsubsection{9.10
検証結果の総括}\label{ux691cux8a3cux7d50ux679cux306eux7dcfux62ec}

{\def\LTcaptype{none} % do not increment counter
\begin{longtable}[]{@{}
  >{\centering\arraybackslash}p{(\linewidth - 10\tabcolsep) * \real{0.1132}}
  >{\centering\arraybackslash}p{(\linewidth - 10\tabcolsep) * \real{0.0755}}
  >{\centering\arraybackslash}p{(\linewidth - 10\tabcolsep) * \real{0.1698}}
  >{\centering\arraybackslash}p{(\linewidth - 10\tabcolsep) * \real{0.3019}}
  >{\centering\arraybackslash}p{(\linewidth - 10\tabcolsep) * \real{0.2264}}
  >{\centering\arraybackslash}p{(\linewidth - 10\tabcolsep) * \real{0.1132}}@{}}
\toprule\noalign{}
\begin{minipage}[b]{\linewidth}\centering
性質
\end{minipage} & \begin{minipage}[b]{\linewidth}\centering
節
\end{minipage} & \begin{minipage}[b]{\linewidth}\centering
ソリトン
\end{minipage} & \begin{minipage}[b]{\linewidth}\centering
\(\sigma\)-依存性
\end{minipage} & \begin{minipage}[b]{\linewidth}\centering
矩形波極限
\end{minipage} & \begin{minipage}[b]{\linewidth}\centering
備考
\end{minipage} \\
\midrule\noalign{}
\endhead
\bottomrule\noalign{}
\endlastfoot
巻き数の定義 & §2.3 & \(\checkmark\) & \textbf{非依存} & \(\checkmark\)
& 境界値のみ \\
巻き数の保存 & §2.3 & \(\checkmark\) & \textbf{非依存} & \(\checkmark\)
& 位相的 \\
キンク/反キンクの存在 & §2.1--2 & \(\checkmark\) & 形状のみ &
\(\checkmark\) & 超関数解 \\
ブリーザーの存在 & §2.4 & \(\checkmark\) & 形状のみ & \(\checkmark\) &
矩形パルス \\
衝突後の形状保存 & §3.1--2 & \(\checkmark\) & \textbf{非依存} &
\(\checkmark\) & \(\sigma\) 保存 \\
衝突後の位相シフト & §3.3 & \(\checkmark\) & \textbf{非依存} &
\(\checkmark\) & \(\alpha\) のみ \\
巻き数保存（衝突） & §3.3 & \(\checkmark\) & \textbf{非依存} &
\(\checkmark\) & 位相的 \\
エネルギー保存 & §4.2 & \(\checkmark\) & 依存 & \(\checkmark\)* &
*各\(2\pi\)回転で自己完結 \\
運動量保存 & §4.2 & \(\checkmark\) & 依存 & \(\checkmark\)* & *同上 \\
Derrick回避 & §5.1--2 & --- & --- & \(\checkmark\) & 閉空間 \(S^1\) \\
パラメータ表現 \((N, \theta_0, \alpha)\) & §8 & \(\checkmark\) &
\textbf{非依存} & \(\checkmark\) & \(\sigma\) は従属 \\
\end{longtable}
}

\textbf{定理 9.1（矩形波極限の構造保存性）.}
sine-Gordon方程式のソリトン解が持つ以下の構造的性質は、矩形波極限
\(\sigma \to 0\) において全て保持される。

\begin{enumerate}
\def\labelenumi{\arabic{enumi}.}
\tightlist
\item
  \textbf{巻き数の定義と保存}（位相的性質：\(\sigma\)-非依存）
\item
  \textbf{衝突後の形状保存と位相シフト}（散乱データの性質：\(\sigma\)-非依存）
\item
  \textbf{有限パラメータ表現}
  \((N, \theta_0, \alpha)\)（モジュライ空間の構造：\(\sigma\)-非依存）
\item
  \textbf{エネルギーと運動量の有限性}（各 \(2\pi\)
  回転の自己相殺構造による。命題 9.6a）
\end{enumerate}

\emph{証明.} (1) は命題9.1--9.2、(2) は命題9.5、(3) は命題8.4と
\(\sigma = \cos\alpha\) の従属性、(4)
は命題9.6aのエネルギー自己相殺構造による。\(\square\)

\textbf{系 9.1.}
矩形波記述はソリトンの「粗い近似」ではない。ソリトンの位相的・散乱的構造を\textbf{完全に継承する厳密な極限}であり、失われるのは広がりパラメータ
\(\sigma\)（内部形状の詳細）のみである。

\textbf{系 9.2（適用限界）.} 本定理の成立条件は、命題 9.6a
で明示した本理論の暗黙の前提------起点 \(\phi = 0\)
の固定、終端の真空値（\(2n\pi\)）への強制、整数巻き数------に依拠する。これらの前提は、無限空間
\((-\infty,+\infty)\) または閉曲線 \(S^1\)
上でのみ成立する。開いた弦（有限区間上の境界条件）への拡張には、半整数巻き数の許容と境界反射の理論（Ghoshal-Zamolodchikov
{[}23{]}）が必要であり、本稿の範囲外である。

\begin{center}\rule{0.5\linewidth}{0.5pt}\end{center}

\subsection{10. 結論}\label{ux7d50ux8ad6}

本稿では、sine-Gordon方程式の数理的構造を概観し、以下を示した。

\begin{enumerate}
\def\labelenumi{\arabic{enumi}.}
\item
  ソリトン解（キンク、反キンク、ブリーザー）は有限個のパラメータ
  \((N, \theta_0, \alpha)\) で\textbf{厳密に}特徴づけられる（§8）。
\item
  ソリトンの持つ全ての構造的性質------巻き数の保存、衝突時の形状保存と位相シフト、有限パラメータ表現------は矩形波極限
  \(\sigma \to 0\) においても保持される（§9）。
\item
  エネルギーの有限性は、各 \(2\pi\)
  回転のエネルギーが自己完結する構造（ポテンシャル
  \(V(\phi) = 1-\cos\phi\) が一周期で \(0 \to 0\)
  に復帰する）に基づく。キンク-反キンク対（\(4\pi\)
  回転）は対消滅可能であり、永続的エネルギーは対消滅できない余りのキンクの寄与のみである。
\item
  したがって、ソリトンから矩形波への移行は情報の損失を伴わない。連続場の無限自由度のうち、物理的に意味のある情報は
  \((N, \theta_0, \alpha)\)
  の有限個のパラメータに集約されており、この集約はソリトンの形状によらず成立する。
\end{enumerate}

\textbf{適用限界：} 上記の結論は、本理論の暗黙の前提（命題
9.6a）の下で成立する。

\begin{itemize}
\tightlist
\item
  \textbf{空間の制約}: 起点の位相が真空値 \(\phi = 0\)
  に固定され、終端の位相が真空値 \(\phi = 2n\pi\) に強制される（巻き数
  \(N\) は整数のみ）。境界反射は存在しない。これらの条件は、無限空間
  \((-\infty, +\infty)\) または閉曲線 \(S^1\)
  上でのみ満たされる。開いた弦（有限区間上の境界条件）への拡張は、半整数巻き数（\(\phi = (2n+1)\pi\),
  ポテンシャル最大値）の許容と境界でのソリトン反射の取り扱いを必要とし、本稿の範囲外である。
\item
  \textbf{衝突の制約}:
  本理論は完全可積分系であり、ソリトン衝突は完全弾性衝突として記述される。運動量の保存則を破らないために、衝突前後でソリトンの速度・形状・エネルギーは個別に保存され、変化するのは位置（位相シフト
  \(\Delta x\)）のみである。これは事実上、ソリトン同士が相互作用しない自由粒子として振る舞うことを意味し、エネルギーの散逸・速度の交換・非弾性チャネル（対消滅→放射）は古典論の範囲では発生しない。
\end{itemize}

矩形波記述はソリトンの「粗い近似」ではなく、位相的・散乱的構造を完全に継承する厳密な極限である。非線形波動方程式の解が有限個の離散パラメータに縮約されるという事実は、連続場の理論と離散構造の間の深い対応を示している。

上記の適用限界------特に、異なる質量の粒子を扱えないこと、および衝突後に速度の再配分が起こらないこと------を克服するモデルを
§11 で提示する。

\begin{center}\rule{0.5\linewidth}{0.5pt}\end{center}

\subsection{11.
完全流体における任意位相の完全弾性衝突モデル}\label{ux5b8cux5168ux6d41ux4f53ux306bux304aux3051ux308bux4efbux610fux4f4dux76f8ux306eux5b8cux5168ux5f3eux6027ux885dux7a81ux30e2ux30c7ux30eb}

§9--10
において、sine-Gordonソリトン理論の適用限界を明らかにした。特に、(i)
全てのソリトンが同一質量を持つこと、(ii)
衝突後に速度の再配分が起こらないこと（透過型衝突）、は物理的な弾性衝突の記述として不十分である。

本節では、これらの制約を持たない代替モデルとして、\textbf{異なる比重（質量）を持つ完全流体の任意位相における完全弾性衝突}を、波動方程式の枠組みで定式化する。

本モデルは、完全流体中の縦波（音波）の波動方程式に基づく。完全流体中の波は本質的に縦波であり、伝播速度は媒質の密度（比重）に依存する（\(c = \sqrt{K/\rho}\)）。異なる比重を持つ2つの波の衝突は、運動量保存とエネルギー保存により完全に記述される。以下では、この物理を最も透明に表現するバネ-質点系として定式化する。

\subsubsection{11.1
モデルの構成}\label{ux30e2ux30c7ux30ebux306eux69cbux6210}

バネ定数 \(k\)、自然長 \(l\)
の完全バネ（散逸ゼロ）2本を直列に接続し、両端を距離 \(2l\)
の壁に固定する。2本のバネの接続点に質量 \(M\) の質点を置く。

\[\text{壁}—[\,k,\; l\,]—M—[\,k,\; l\,]—\text{壁}\]

平衡位置（中央）からの変位を \(x\) とすると、左バネの復元力 \(-kx\)
と右バネの復元力 \(-kx\) の合力により、運動方程式は

\[M\ddot{x} + 2kx = 0 \tag{W1}\]

これは角振動数 \(\omega = \sqrt{2k/M}\) の単振動（縦波）である。一般解は

\[x(t) = A\cos(\omega t + \phi_0) \tag{W2}\]

であり、\(A\) は振幅、\(\phi_0\) は初期位相である。

\subsubsection{11.2
2体系の配置と衝突の定義}\label{ux4f53ux7cfbux306eux914dux7f6eux3068ux885dux7a81ux306eux5b9aux7fa9}

異なる質量 \(M_1, M_2\) を持つ2つの独立な単振動系を考える。

\[\ddot{x}_1 + \omega_1^2 x_1 = 0, \qquad \omega_1 = \sqrt{\frac{2k}{M_1}} \tag{W3}\]

\[\ddot{x}_2 + \omega_2^2 x_2 = 0, \qquad \omega_2 = \sqrt{\frac{2k}{M_2}} \tag{W4}\]

\(M_1 \neq M_2\) のとき \(\omega_1 \neq \omega_2\)
であり、2つの系は独立に振動する。完全流体の比重差は、この質量差
\(M_1 \neq M_2\) に対応する。

\textbf{衝突の定義.} 本稿における「衝突」とは、両系の状態
\((x_1, \dot{x}_1)\) と \((x_2, \dot{x}_2)\) が、ある時刻 \(t = t_c\)
において運動量保存則とエネルギー保存則を介して相互作用し、瞬間的に新しい状態
\((x_1, \dot{x}_1')\) と \((x_2, \dot{x}_2')\)
に遷移する数学的過程として定義する。位置 \(x_i\)
は衝突によって変化せず、速度 \(v_i = \dot{x}_i\) のみが (W7)--(W8)
に従って変化する。

\textbf{注記.}
本モデルにおいて、衝突は物理的接触の幾何学的記述ではなく、「2つの振動系の状態が保存則を介して結合する瞬間的な状態変換」として抽象化されている。完全流体中の縦波の衝突における物理的接触面の役割を、保存則による状態変換に置き換えたものである。衝突の発生時刻
\(t_c\)
は外部から与えられる入力パラメータであり、本モデルの内部からは決定されない。これにより、任意の位相
\(\phi_1, \phi_2\) における衝突を統一的に扱うことができる。

\subsubsection{11.3
任意位相での完全弾性衝突}\label{ux4efbux610fux4f4dux76f8ux3067ux306eux5b8cux5168ux5f3eux6027ux885dux7a81}

\textbf{衝突直前の状態：} 時刻 \(t_c\) における各系の状態を以下で表す。

\[x_{1,c} = A_1\cos\phi_1, \qquad v_1 = -A_1\omega_1\sin\phi_1 \tag{W5}\]

\[x_{2,c} = A_2\cos\phi_2, \qquad v_2 = -A_2\omega_2\sin\phi_2 \tag{W6}\]

ここで \(\phi_i = \omega_i t_c + \phi_{0,i}\)
は衝突瞬間における各系の位相、\(A_i\) は振幅である。速度の符号は
(W3)(W4) の自然な微分から定まる。

\textbf{完全弾性衝突の条件（運動量保存 + エネルギー保存）：}

\[M_1 v_1 + M_2 v_2 = M_1 v_1' + M_2 v_2' \tag{W7}\]

\[\frac{1}{2}M_1 v_1^2 + \frac{1}{2}M_2 v_2^2 = \frac{1}{2}M_1 v_1'^2 + \frac{1}{2}M_2 v_2'^2 \tag{W8}\]

この2条件から衝突後の速度が一意に決定される。

\textbf{定理 11.1（任意位相の弾性衝突公式）.} 位相 \(\phi_1, \phi_2\)
で衝突する質量 \(M_1, M_2\) の完全弾性衝突後の速度は以下で与えられる。

\[v_1' = \frac{(M_1 - M_2)\,v_1 + 2M_2\,v_2}{M_1 + M_2} \tag{W9}\]

\[v_2' = \frac{(M_2 - M_1)\,v_2 + 2M_1\,v_1}{M_1 + M_2} \tag{W10}\]

\emph{証明.} (W7) と (W8)
の連立方程式を解くことにより直ちに得られる。\(\square\)

\subsubsection{11.4
衝突後の振動状態}\label{ux885dux7a81ux5f8cux306eux632fux52d5ux72b6ux614b}

衝突は瞬間的であり、位置は変化しない。衝突後、各質点は衝突時の位置
\(x_{i,c}\) と新速度 \(v_i'\)
を初期条件として、元のバネ系で振動を再開する。

\[x_1(t) = x_{1,c}\cos(\omega_1\tau) + \frac{v_1'}{\omega_1}\sin(\omega_1\tau) \qquad (\tau = t - t_c) \tag{W11}\]

\[x_2(t) = x_{2,c}\cos(\omega_2\tau) + \frac{v_2'}{\omega_2}\sin(\omega_2\tau) \tag{W12}\]

\textbf{衝突後の振幅と位相：}

\[A_i' = \sqrt{x_{i,c}^2 + \left(\frac{v_i'}{\omega_i}\right)^2} \tag{W13}\]

\[\phi_i' = \arctan\!\left(\frac{-v_i'}{\omega_i\, x_{i,c}}\right) \tag{W14}\]

角振動数 \(\omega_i\)
はバネ定数と質量のみで決まるため、衝突によって変化しない。変化するのは振幅
\(A_i\) と位相 \(\phi_i\) のみである。

\subsubsection{11.5
エネルギー保存の証明}\label{ux30a8ux30cdux30ebux30aeux30fcux4fddux5b58ux306eux8a3cux660e}

\textbf{命題 11.1.}
任意位相の完全弾性衝突において、系の全エネルギーは保存される。

\emph{証明.}
各系の全エネルギーはバネのポテンシャルエネルギーと運動エネルギーの和である。

\[E_i = \frac{1}{2}M_i\omega_i^2\,x_{i,c}^2 + \frac{1}{2}M_i\,v_i^2 = k\,x_{i,c}^2 + \frac{1}{2}M_i\,v_i^2\]

衝突は瞬間的であるため、位置エネルギー \(k\,x_{i,c}^2\)
は衝突前後で不変である。運動エネルギーは弾性衝突の条件 (W8)
により保存される。したがって

\[E_1 + E_2 = E_1' + E_2' \qquad \square \tag{W15}\]

\subsubsection{11.6
衝突位相による振幅・位相変化の特徴}\label{ux885dux7a81ux4f4dux76f8ux306bux3088ux308bux632fux5e45ux4f4dux76f8ux5909ux5316ux306eux7279ux5fb4}

(W9)--(W10) の弾性衝突公式と (W13)(W14)
による衝突後振幅・位相の決定式から、衝突瞬間の位相 \(\phi_i\)
に応じて以下の特徴が現れる。

{\def\LTcaptype{none} % do not increment counter
\begin{longtable}[]{@{}
  >{\raggedright\arraybackslash}p{(\linewidth - 8\tabcolsep) * \real{0.2000}}
  >{\raggedright\arraybackslash}p{(\linewidth - 8\tabcolsep) * \real{0.2000}}
  >{\raggedright\arraybackslash}p{(\linewidth - 8\tabcolsep) * \real{0.2000}}
  >{\raggedright\arraybackslash}p{(\linewidth - 8\tabcolsep) * \real{0.2000}}
  >{\raggedright\arraybackslash}p{(\linewidth - 8\tabcolsep) * \real{0.2000}}@{}}
\toprule\noalign{}
\begin{minipage}[b]{\linewidth}\raggedright
衝突位相
\end{minipage} & \begin{minipage}[b]{\linewidth}\raggedright
位置 \(x_{i,c}\)
\end{minipage} & \begin{minipage}[b]{\linewidth}\raggedright
速度 \(|v_i|\)
\end{minipage} & \begin{minipage}[b]{\linewidth}\raggedright
エネルギー配分
\end{minipage} & \begin{minipage}[b]{\linewidth}\raggedright
衝突の効果
\end{minipage} \\
\midrule\noalign{}
\endhead
\bottomrule\noalign{}
\endlastfoot
\(\phi_i = \pi/2\) & \(0\) & \(A_i\omega_i\)（最大） &
すべて運動エネルギー & 速度が完全に再配分され、振幅 \(A_i\)
のみが変化する \\
\(\phi_i = 0\) & \(A_i\)（最大） & \(0\) & すべて位置エネルギー &
速度の交換が発生せず、衝突前後で状態は変化しない \\
\(\phi_i = \pi/4\) & \(A_i/\sqrt{2}\) & \(A_i\omega_i/\sqrt{2}\) &
位置と運動が等分配 & 速度の部分的再配分と位相のずれが発生する \\
\end{longtable}
}

\textbf{注記.} \(\phi_i = 0\)
において「状態が変化しない」のは、(W9)(W10) の右辺が \(v_i\)
のみで定まり、\(v_i = 0\) のとき \(v_i' = 0\)
が解となるためである。これは衝突が発生しないことを意味するのではなく、速度ゼロの状態同士が保存則を介して結合した結果、状態が不変であったことを意味する。一般的な
\(\phi_1 \neq \phi_2\)
の場合、両系の速度が同時にゼロでない限り、必ず非自明な状態変換が発生する。

\subsubsection{11.7
sine-Gordon理論との比較}\label{sine-gordonux7406ux8ad6ux3068ux306eux6bd4ux8f03}

{\def\LTcaptype{none} % do not increment counter
\begin{longtable}[]{@{}
  >{\raggedright\arraybackslash}p{(\linewidth - 4\tabcolsep) * \real{0.3333}}
  >{\raggedright\arraybackslash}p{(\linewidth - 4\tabcolsep) * \real{0.3333}}
  >{\raggedright\arraybackslash}p{(\linewidth - 4\tabcolsep) * \real{0.3333}}@{}}
\toprule\noalign{}
\begin{minipage}[b]{\linewidth}\raggedright
性質
\end{minipage} & \begin{minipage}[b]{\linewidth}\raggedright
sine-Gordon（§1--10）
\end{minipage} & \begin{minipage}[b]{\linewidth}\raggedright
本モデル（§11）
\end{minipage} \\
\midrule\noalign{}
\endhead
\bottomrule\noalign{}
\endlastfoot
異なる質量 & 不可（全キンク同一質量） &
\textbf{可}（\(M_1 \neq M_2\)） \\
衝突後の速度再配分 & なし（透過型） & \textbf{あり}（弾性衝突公式） \\
任意位相での衝突 & 未定義 & \textbf{完全に定義} \\
閉じた弦の仮定 & 必要（暗黙） & \textbf{不要} \\
巻き数の整数制約 & 必要 & \textbf{不要} \\
半整数位相 & 排除（\(\phi = \pi\) 不許容） & \textbf{許容} \\
波動方程式 & \(\phi_{tt} - \phi_{xx} + \sin\phi = 0\) &
\(\ddot{x}_i + \omega_i^2 x_i = 0\) \\
隠れた仮定 & 6つ（命題 9.6a） & \textbf{なし}（保存則のみ） \\
エネルギー散逸 & ゼロ & ゼロ \\
定式化の前提 & 運動量保存 + エネルギー保存 + 位相的制約 + 可積分性 +
閉空間 + 同一質量 & \textbf{運動量保存 + エネルギー保存のみ} \\
\end{longtable}
}

\subsubsection{11.8 物理的解釈}\label{ux7269ux7406ux7684ux89e3ux91c8}

本モデルにおけるパラメータの物理的対応を以下に示す。

{\def\LTcaptype{none} % do not increment counter
\begin{longtable}[]{@{}ll@{}}
\toprule\noalign{}
モデルのパラメータ & 物理的対応 \\
\midrule\noalign{}
\endhead
\bottomrule\noalign{}
\endlastfoot
質量 \(M_i\) & 完全流体の比重（密度） \\
振幅 \(A_i\) & 波の振幅（エネルギーの大きさ） \\
位相 \(\phi_i\) & 波の位相 \\
角振動数 \(\omega_i = \sqrt{2k/M_i}\) & 固有振動数（比重に依存） \\
衝突（状態変換） & 保存則を介した波の瞬間的相互作用 \\
\end{longtable}
}

完全流体中の縦波が完全弾性衝突する過程は、(W1)--(W15)
により完全に記述される。仮定は運動量保存とエネルギー保存の2つの保存則のみであり、位相的制約、可積分性、空間のコンパクト性、巻き数の整数性などの追加仮定は一切不要である。

\subsubsection{11.9
矩形波極限との同値性}\label{ux77e9ux5f62ux6ce2ux6975ux9650ux3068ux306eux540cux5024ux6027}

§9
において、sine-Gordon方程式のソリトンが矩形波極限でも構造を保存することを示した。本節では、§11
のモデルにおいても、波形が正弦波（余弦波）から矩形波に変化しても衝突の結果が同一であることを証明する。

\textbf{波形パラメータの導入.} 振動の波形を、遷移幅パラメータ
\(\sigma\)（\(0 < \sigma \leq 1\)）で制御する族 \(x^\sigma(t)\)
として定義する。

\begin{itemize}
\tightlist
\item
  \(\sigma = 1\)：滑らかな正弦波（余弦波）\(x^1(t) = A\cos(\omega t)\)
\item
  \(\sigma \to 0\)：矩形波 \(x^0(t) = A\,\mathrm{sgn}(\cos(\omega t))\)
\end{itemize}

中間的な \(\sigma\) は、遷移領域の急峻さを制御する。全ての \(\sigma\)
に対して、振動の振幅は \(A\)、周期は \(T = 2\pi/\omega\) で共通である。

\textbf{定理 11.2（衝突結果の波形非依存性）.} 衝突時の位置 \(x_{i,c}\)
と全エネルギー \(E_i\) が与えられたとき、衝突後の速度 \(v_i'\)、振幅
\(A_i'\)、位相 \(\phi_i'\) は波形パラメータ \(\sigma\) に依存しない。

\emph{証明.} 3段階で示す。

\textbf{第1段階：衝突時の速度はエネルギーと位置から一意に決まる。}

バネ-質点系の全エネルギーは、波形によらず保存される。

\[E_i = k\,x_i^2(t) + \frac{1}{2}M_i\,\dot{x}_i^2(t) = k\,A_i^2 \tag{W16}\]

これは波形に依存しない。バネのポテンシャルエネルギーは位置のみの関数であり、運動エネルギーは残りである。したがって、衝突位置
\(x_{i,c}\) における速度は

\[v_i^2 = \frac{2(E_i - k\,x_{i,c}^2)}{M_i} = \omega_i^2(A_i^2 - x_{i,c}^2) \tag{W17}\]

で\textbf{一意に}決定される。これは \(\sigma\)
に依存しない。正弦波でも矩形波でも、同じ位置を同じエネルギーで通過する質点は、同じ速度を持つ。\(\checkmark\)

\textbf{第2段階：弾性衝突公式は速度のみに依存する。}

定理 11.1 の衝突公式 (W9)--(W10) は、\(M_1, M_2, v_1, v_2\)
のみの代数的関係式であり、波形の情報を一切含まない。第1段階で \(v_i\) が
\(\sigma\)-非依存であることが示されたので、\(v_i'\) も
\(\sigma\)-非依存である。\(\checkmark\)

\textbf{第3段階：衝突後の振幅と位相も \(\sigma\)-非依存である。}

衝突後の全エネルギーは

\[E_i' = k\,x_{i,c}^2 + \frac{1}{2}M_i\,v_i'^2 \tag{W18}\]

であり、新しい振幅は

\[A_i' = \sqrt{\frac{E_i'}{k}} = \sqrt{x_{i,c}^2 + \frac{v_i'^2}{\omega_i^2}} \tag{W19}\]

\(x_{i,c}\)、\(v_i'\)、\(\omega_i\) の全てが
\(\sigma\)-非依存であるから、\(A_i'\) も \(\sigma\)-非依存である。

初期位相 \(\phi_i' = \arctan(-v_i'/(\omega_i\,x_{i,c}))\) も同様に
\(\sigma\)-非依存である。\(\square\)

\textbf{系 11.1（矩形波極限の厳密性）.}
本モデルにおける完全弾性衝突の結果は、振動が滑らかな正弦波であっても矩形波であっても\textbf{厳密に同一}である。失われる情報は波形の詳細（遷移領域の形状）のみであり、衝突の物理------速度の再配分、エネルギー保存、振幅と位相の変化------は波形に一切依存しない。

\textbf{系 11.2.} 本結果は、§9 で
sine-Gordon方程式について示した矩形波極限の構造保存性を、より一般的な枠組みで再現するものである。sine-Gordonでは構造保存性の証明に可積分性と閉空間の仮定が必要であったが、本モデルでは\textbf{エネルギー保存の1条件のみ}（式
(W16)）で同値性が成立する。

\subsubsection{11.10
本モデルの検証総括}\label{ux672cux30e2ux30c7ux30ebux306eux691cux8a3cux7dcfux62ec}

{\def\LTcaptype{none} % do not increment counter
\begin{longtable}[]{@{}
  >{\centering\arraybackslash}p{(\linewidth - 8\tabcolsep) * \real{0.1364}}
  >{\centering\arraybackslash}p{(\linewidth - 8\tabcolsep) * \real{0.1818}}
  >{\centering\arraybackslash}p{(\linewidth - 8\tabcolsep) * \real{0.1818}}
  >{\centering\arraybackslash}p{(\linewidth - 8\tabcolsep) * \real{0.3636}}
  >{\centering\arraybackslash}p{(\linewidth - 8\tabcolsep) * \real{0.1364}}@{}}
\toprule\noalign{}
\begin{minipage}[b]{\linewidth}\centering
性質
\end{minipage} & \begin{minipage}[b]{\linewidth}\centering
正弦波
\end{minipage} & \begin{minipage}[b]{\linewidth}\centering
矩形波
\end{minipage} & \begin{minipage}[b]{\linewidth}\centering
\(\sigma\)-依存性
\end{minipage} & \begin{minipage}[b]{\linewidth}\centering
備考
\end{minipage} \\
\midrule\noalign{}
\endhead
\bottomrule\noalign{}
\endlastfoot
衝突時の速度 \(v_i\) & \(\checkmark\) & \(\checkmark\) & \textbf{非依存}
& (W17) による \\
衝突後の速度 \(v_i'\) & \(\checkmark\) & \(\checkmark\) &
\textbf{非依存} & (W9)--(W10) \\
衝突後の振幅 \(A_i'\) & \(\checkmark\) & \(\checkmark\) &
\textbf{非依存} & (W19) \\
衝突後の位相 \(\phi_i'\) & \(\checkmark\) & \(\checkmark\) &
\textbf{非依存} & (W14) \\
エネルギー保存 & \(\checkmark\) & \(\checkmark\) & \textbf{非依存} &
(W15) \\
運動量保存 & \(\checkmark\) & \(\checkmark\) & \textbf{非依存} & (W7) \\
異なる質量 & \(\checkmark\) & \(\checkmark\) & \textbf{非依存} &
\(M_1 \neq M_2\) \\
任意位相 & \(\checkmark\) & \(\checkmark\) & \textbf{非依存} &
\(\phi_1, \phi_2\) 任意 \\
\end{longtable}
}

\textbf{定理 11.3（完全弾性衝突モデルの構造保存性）.} 異なる比重（質量
\(M_1 \neq M_2\)）を持つ完全流体の縦波が任意位相 \(\phi_1, \phi_2\)
で完全弾性衝突する過程は、波動方程式 (W3)--(W4) と弾性衝突公式
(W9)--(W10) により完全に記述され、以下の全ての性質が矩形波極限
\(\sigma \to 0\) において保持される。

\begin{enumerate}
\def\labelenumi{\arabic{enumi}.}
\tightlist
\item
  \textbf{速度の再配分}（弾性衝突公式：\(\sigma\)-非依存）
\item
  \textbf{エネルギーと運動量の保存}（(W15) と (W7)：\(\sigma\)-非依存）
\item
  \textbf{振幅と位相の決定}（(W19) と (W14)：\(\sigma\)-非依存）
\end{enumerate}

\emph{証明.} 定理 11.2 の直接の帰結である。\(\square\)

本定式化の全仮定は以下の2つのみである。

\begin{itemize}
\tightlist
\item
  \textbf{運動量保存}: \(M_1 v_1 + M_2 v_2 = M_1 v_1' + M_2 v_2'\)
\item
  \textbf{エネルギー保存}:
  \(\frac{1}{2}M_1 v_1^2 + \frac{1}{2}M_2 v_2^2 = \frac{1}{2}M_1 v_1'^2 + \frac{1}{2}M_2 v_2'^2\)
\end{itemize}

位相的制約、可積分性、空間のコンパクト性、巻き数の整数性、ポテンシャルの周期性、同一質量の仮定は一切不要であり、矩形波極限との同値性はエネルギー保存から自動的に導かれる。

\begin{center}\rule{0.5\linewidth}{0.5pt}\end{center}

\subsection{参考文献}\label{ux53c2ux8003ux6587ux732e}

{[}1{]} Frenkel, J. \& Kontorova, T. ``On the theory of plastic
deformation and twinning.'' \emph{Izv. Akad. Nauk SSSR, Ser. Fiz.} 1,
137--149 (1939).

{[}2{]} Scott, A.C., Chu, F.Y.F. \& McLaughlin, D.W. ``The Soliton: A
New Concept in Applied Science.'' \emph{Proc. IEEE} 61, 1443--1483
(1973).

{[}3{]} Rubinstein, J. ``Sine-Gordon equation.'' \emph{J. Math. Phys.}
11, 258--266 (1970).

{[}4{]} Rajaraman, R. \emph{Solitons and Instantons: An Introduction to
Solitons and Instantons in Quantum Field Theory.} North-Holland (1982).

{[}5{]} Bour, E. ``Théorie de la déformation des surfaces.'' \emph{J.
Éc. Imp. Polytech.} 22, Cahier 39, 1--148 (1862).

{[}6{]} Enneper, A. ``Über asymptotische Linien.'' \emph{Nachr. Königl.
Gesellsch. d.~Wiss. Göttingen} 1870, 493--510 (1870).

{[}7{]} Bianchi, L. ``Ricerche sulle superficie a curvatura costante e
sulle elicoidi.'' \emph{Ann. Sc. Norm. Super. Pisa} 2, 285--341 (1879).

{[}8{]} Skyrme, T.H.R. ``A non-linear field theory.'' \emph{Proc. R.
Soc. Lond. A} 260, 127--138 (1961).

{[}9{]} Perring, J.K. \& Skyrme, T.H.R. ``A model unified field
equation.'' \emph{Nucl. Phys.} 31, 550--555 (1962).

{[}10{]} Manton, N.S. \& Sutcliffe, P.M. \emph{Topological Solitons.}
Cambridge University Press (2004).

{[}11{]} Ablowitz, M.J., Kaup, D.J., Newell, A.C. \& Segur, H. ``Method
for solving the sine-Gordon equation.'' \emph{Phys. Rev.~Lett.} 30,
1262--1264 (1973).

{[}12{]} Ablowitz, M.J., Kaup, D.J., Newell, A.C. \& Segur, H. ``The
inverse scattering transform --- Fourier analysis for nonlinear
problems.'' \emph{Stud. Appl. Math.} 53, 249--315 (1974).

{[}13{]} Faddeev, L.D. \& Takhtajan, L.A. \emph{Hamiltonian Methods in
the Theory of Solitons.} Springer (1987).

{[}14{]} Zakharov, V.E. \& Shabat, A.B. ``Exact theory of
two-dimensional self-focusing and one-dimensional self-modulation of
waves in nonlinear media.'' \emph{Sov. Phys. JETP} 34, 62--69 (1972).

{[}15{]} Rogers, C. \& Schief, W.K. \emph{Bäcklund and Darboux
Transformations: Geometry and Modern Applications in Soliton Theory.}
Cambridge University Press (2002).

{[}16{]} Gibbon, J.D., James, I.N. \& Moroz, I.M. ``The sine-Gordon
equation as a model for a rapidly rotating baroclinic fluid.''
\emph{Proc. R. Soc. Lond. A} 367, 219--237 (1979).

{[}17{]} Derrick, G.H. ``Comments on nonlinear wave equations as models
for elementary particles.'' \emph{J. Math. Phys.} 5, 1252--1254 (1964).

{[}18{]} Manton, N.S. ``Geometry of Skyrmions.'' \emph{Commun. Math.
Phys.} 111, 469--478 (1987).

{[}19{]} Barone, A. \& Paternò, G. \emph{Physics and Applications of the
Josephson Effect.} Wiley (1982).

{[}20{]} Coleman, S. ``Quantum sine-Gordon equation as the massive
Thirring model.'' \emph{Phys. Rev.~D} 11, 2088--2097 (1975).

{[}21{]} Mandelstam, S. ``Soliton operators for the quantized
sine-Gordon equation.'' \emph{Phys. Rev.~D} 11, 3026--3030 (1975).

{[}22{]} Manton, N.S. ``A remark on the scattering of BPS monopoles.''
\emph{Phys. Lett. B} 110, 54--56 (1982).

{[}23{]} Ghoshal, S. \& Zamolodchikov, A.B. ``Boundary S matrix and
boundary state in two-dimensional integrable quantum field theory.''
\emph{Int. J. Mod. Phys. A} 9, 3841--3885 (1994).

\end{document}
